\documentclass{article}
\usepackage[T2A]{fontenc}
\usepackage[utf8]{inputenc}
\usepackage[russian]{babel}
\usepackage{amsmath, amssymb, amsthm}
\usepackage{amsfonts}
\usepackage{dsfont}
\usepackage{hyperref}
\usepackage{enumitem}
\usepackage{geometry}
\geometry{a4paper, margin=2.5cm}

\newtheorem{theorem}{Теорема}
\newtheorem{axiom}{Аксиома}

\title{Математическая формализация манифеста устройства \\
    в системе <<Умный огород>>}
\author{}
\date{}

\begin{document}

\maketitle


\section{Базовые определения}

\subsection{Пространство типов}
Определим множество примитивных типов данных:

\[
    \mathbb{T} = \{\text{String}, \text{Integer}, \text{Float}, \text{Boolean}, \text{Timestamp}, \text{Enum}, \text{Array}, \text{Object}\}
\]

\subsection{Пространство значений}
Каждому типу $t \in \mathbb{T}$ соответствует множество допустимых значений $V_t$:

\begin{itemize}
    \item $V_{\text{String}} = \Sigma^*$, где $\Sigma$ — конечный алфавит
    \item $V_{\text{Integer}} = \mathbb{Z}$
    \item $V_{\text{Float}} = \mathbb{R}$
    \item $V_{\text{Boolean}} = \{\text{true}, \text{false}\}$
    \item $V_{\text{Timestamp}} = \mathbb{R}_{\geq 0}$ (Unix time)
    \item $V_{\text{Enum}} = \{e_1, e_2, \ldots, e_n\}$ (конечное множество)
    \item $V_{\text{Array}} = \bigcup_{n=0}^{\infty} V^n$ (множество всех кортежей)
    \item $V_{\text{Object}} = \mathcal{P}(\text{Key} \times V)$ (частичные функции из ключей в значения)
\end{itemize}

\section{Формальное определение манифеста}

\subsection{Манифест как кортеж}
Манифест устройства $M$ определим как упорядоченный кортеж:

\[
M = (V, I, T, P, C, N, F, U, S) \in \mathcal{M}
\]

где $\mathcal{M}$ — пространство всех допустимых манифестов.

\subsection{Компоненты манифеста}

\subsubsection{Версия манифеста}
\[
V \in \mathbb{N}^2 \times \mathbb{N}
\]

$V = (maj, min, patch)$, где $maj, min, patch \in \mathbb{N}$ — семантическая версия.

\subsubsection{Информация об устройстве}
\[
I = (id, man, model, rev) \in \mathcal{I}
\]

где:
\begin{itemize}
    \item $id \in \text{String}$ — уникальный идентификатор в формате UUID
    \item $man \in \text{String}$ — производитель
    \item $model \in \text{String}$ — модель
    \item $rev \in \text{String}$ — ревизия
\end{itemize}

При этом существует инъективная функция идентификации:
\[
\phi: \mathcal{I} \to \text{String}
\]
\[
\phi(I) = id
\]

\subsubsection{Тип устройства}
\[
T \in \mathcal{T} = \{\text{sensor}, \text{actuator}, \text{controller}, \text{gateway}\}
\]

Дополнительно определим функцию категоризации:
\[
\kappa: \mathcal{T} \to \mathcal{C}
\]
где $\mathcal{C}$ — множество категорий устройств (например, soil\_moisture\_sensor, water\_valve и т.д.)

\subsubsection{Протоколы взаимодействия}
\[
P = (p_{primary}, \mathcal{P}_{backup}, v_p) \in \mathcal{P}
\]

где:
\begin{itemize}
    \item $p_{primary} \in \Pi$ — основной протокол из множества протоколов $\Pi$
    \item $\mathcal{P}_{backup} \subset \Pi$ — множество резервных протоколов
    \item $v_p \in \text{String}$ — версия протокола
\end{itemize}

\subsubsection{Возможности устройства}
\[
C = (M_{sens}, A_{ctrl}, F_{meas}) \in \mathcal{C}
\]

где:
\begin{itemize}
    \item $M_{sens} \subset \mathcal{S}$ — множество сенсоров
    \item $A_{ctrl} \subset \mathcal{A}$ — множество актуаторов
    \item $F_{meas} \subset \mathcal{F}$ — множество измерительных функций
\end{itemize}

Сенсор определяется как:
\[
s = (type, unit, range, accuracy, interval) \in \mathcal{S}
\]
\[
range = (min, max) \in \mathbb{R}^2
\]
\[
accuracy \in \mathbb{R}_{\geq 0}
\]
\[
interval \in \mathbb{R}_{>0} \cup \{\text{on-demand}\}
\]

Актуатор определяется как:
\[
a = (type, states, default) \in \mathcal{A}
\]
\[
states \subset \mathbb{S} \text{ (конечное множество состояний)}
\]
\[
default \in states
\]

\subsubsection{Сетевые параметры}
\[
N = (m, \pi, f, q) \in \mathcal{N}
\]

где:
\begin{itemize}
    \item $m \in \{\text{true}, \text{false}\}$ — поддержка mesh-сети
    \item $\pi \in \text{String}$ — mesh-протокол
    \item $f \in \mathbb{R}_{>0}$ — рабочая частота в ГГц
    \item $q \in [0, 1] \times \mathbb{R}$ — качество связи $(link\_quality, rssi)$
\end{itemize}

\subsubsection{Прошивка}
\[
F = (v_f, d_f, ota) \in \mathcal{F}
\]

где:
\begin{itemize}
    \item $v_f \in \text{String}$ — версия прошивки
    \item $d_f \in \text{Timestamp}$ — дата сборки
    \item $ota \in \text{Boolean}$ — поддержка обновлений "по воздуху"
\end{itemize}

\subsubsection{UI-подсказки}
\[
U = (icon, color, desc) \in \mathcal{U}
\]

где:
\begin{itemize}
    \item $icon \in \text{String}$ — идентификатор иконки
    \item $color \in \text{String}$ — цвет в HEX-формате
    \item $desc \in \text{String}$ — локализованное описание
\end{itemize}

\subsubsection{Совместимые сценарии}
\[
S \subset \Theta
\]

где $\Theta$ — множество всех доступных сценариев использования.

\section{Формальные отношения и ограничения}

\subsection{Отношение совместимости}
Определим бинарное отношение совместимости между устройством и сценарием:

\[
\vDash \subseteq \mathcal{M} \times \Theta
\]

$M \vDash \theta$ означает, что устройство $M$ совместимо со сценарием $\theta$.

\subsection{Аксиомы корректности}

Для любого манифеста $M$ должны выполняться следующие аксиомы:

\begin{axiom}[Идентификации]
\[
\forall M_1, M_2 \in \mathcal{M}: \phi(I_1) = \phi(I_2) \implies M_1 = M_2
\]
\end{axiom}

\begin{axiom}[Полноты]
\[
\forall M \in \mathcal{M}: (T = \text{sensor}) \implies (M_{sens} \neq \varnothing)
\]
\[
\forall M \in \mathcal{M}: (T = \text{actuator}) \implies (A_{ctrl} \neq \varnothing)
\]
\end{axiom}

\begin{axiom}[Согласованности]
\[
\forall s \in M_{sens}: min \leq max
\]
\[
\forall s \in M_{sens}: interval \in \mathbb{R}_{>0} \implies interval \geq \epsilon \text{ (минимальный интервал)}
\]
\end{axiom}

\subsection{Функция валидации}

Определим функцию валидации манифеста:

\[
\text{valid}: \mathcal{M} \to \{\text{true}, \text{false}\}
\]

\[
\text{valid}(M) = \bigwedge_{i=1}^{n} \text{constraint}_i(M)
\]

где $\text{constraint}_i$ — проверка выполнения i-го ограничения.

\section{Операции над манифестами}

\subsection{Сравнение версий}
\[
\text{version\_cmp}: \mathbb{N}^3 \times \mathbb{N}^3 \to \{-1, 0, 1\}
\]

\[
\text{version\_cmp}((a_1,a_2,a_3), (b_1,b_2,b_3)) = 
\begin{cases}
-1 & \text{если } a_1 < b_1 \text{ или } (a_1 = b_1 \land a_2 < b_2) \text{ или } (a_1 = b_1 \land a_2 = b_2 \land a_3 < b_3) \\
0 & \text{если } a_1 = b_1 \land a_2 = b_2 \land a_3 = b_3 \\
1 & \text{иначе}
\end{cases}
\]

\subsection{Слияние манифестов (для составных устройств)}
Для устройств с дочерними элементами определим операцию слияния:

\[
\oplus: \mathcal{M} \times \mathcal{M}^* \to \mathcal{M}
\]

\[
M \oplus [M_1, M_2, \ldots, M_n] = M'
\]

где $M'$ содержит информацию из $M$ и всех дочерних манифестов.

\section{Метрика расстояния между манифестами}

Для поиска похожих устройств определим метрику:

\[
d: \mathcal{M} \times \mathcal{M} \to \mathbb{R}_{\geq 0}
\]

\[
d(M_1, M_2) = \sum_{i} w_i \cdot \delta_i(M_1, M_2)
\]

где $w_i$ — веса, а $\delta_i$ — функции различия по каждому компоненту.

\section{Формальная грамматика манифеста}

В нотации расширенной формы Бэкуса-Наура (EBNF):

\begin{verbatim}
manifest   = version , device-info , type , protocols , capabilities , 
              network , firmware , ui-hints , scenarios ;

version    = major , "." , minor , "." , patch ;
major      = number ;
minor      = number ;
patch      = number ;

device-info = "{" , id , "," , manufacturer , "," , model , "," , revision , "}" ;
id         = string , ":" , hexdigit{8} , "-" , hexdigit{4} , "-" , hexdigit{4} , 
             "-" , hexdigit{4} , "-" , hexdigit{12} ;

capabilities = sensors , actuators ;
sensors    = "[" , sensor { "," , sensor } , "]" ;
sensor     = "{" , type , "," , unit , "," , range , "," , accuracy , "," , interval , "}" ;
range      = "(" , number , "," , number , ")" ;
\end{verbatim}

\section{Теорема о единственности манифеста}

\begin{theorem}[О единственности манифеста]
Для каждого физического устройства существует единственный допустимый манифест в момент времени $t$.
\end{theorem}

\begin{proof}
Следует из аксиомы идентификации и детерминированности процесса генерации манифеста. Если предположить существование двух различных манифестов $M_1 \neq M_2$ для одного устройства, то $\phi(I_1) = \phi(I_2)$ (один и тот же ID), что противоречит аксиоме идентификации.
\end{proof}

\paragraph{Итог:} Данная математическая формализация обеспечивает строгую основу для реализации системы самоописания устройств, гарантируя однозначность интерпретации, возможность автоматической валидации и совместимость компонентов экосистемы.

\end{document}